\documentclass[10pt]{article}

% Specify the margins and text size.
\setlength{\textwidth}{6.5in}
\setlength{\textheight}{9.6in}
\setlength{\oddsidemargin}{0pt}
\setlength{\evensidemargin}{0pt}
\setlength{\topmargin}{0pt}
\setlength{\hoffset}{-.05in}
\setlength{\voffset}{-0.8in}

\setlength{\parskip}{5pt}
\setlength{\parindent}{0pt}

% Load some fonts and symbol packages
\usepackage{latexsym}
\usepackage{pifont}       % contains 'star' symbol for counterinsurgency handbook title
%\usepackage{yfonts} 
\usepackage{amsmath}
\usepackage{amsfonts}

\usepackage{graphicx}     % actually, this is loaded by pstricks
\usepackage[T1]{fontenc}
\usepackage{ifthen}
\usepackage{pstricks,pst-grad,pst-text,pst-node,multido,pst-plot,calc,pst-3dplot}
%\usepackage[all]{xy}
%\usepackage{animate}

% The hyperref package inserts links into the pdf.
\definecolor{MyLinkColor}{rgb}{.1,.2,1}
\definecolor{MyCiteColor}{rgb}{.1,1,.2}
\definecolor{MyURLColor}{rgb}{.4,.4,.4}
\usepackage[backref=true,pagebackref=false,hyperindex,colorlinks=true,
  linkcolor=MyLinkColor,urlcolor=MyURLColor]{hyperref}


% The tweaklist package is something I found on the web.  It provides a simple interface
% for making changes to spacing used in the itemize and enumerate environments.  Comment
% this out if you don't care to use tweaklist.
\usepackage{tweaklist}
\renewcommand{\itemhook}{\setlength{\parskip}{2pt}\setlength{\parsep}%
{1pt}\setlength{\topsep}{0pt}\setlength{\itemsep}{0pt}}
\renewcommand{\enumhook}{\setlength{\parskip}{2pt}\setlength{\parsep}%
{1pt}\setlength{\topsep}{0pt}\setlength{\itemsep}{0pt}}

\begin{document}
\pagestyle{plain}
\setlength{\baselineskip}{1.05\baselineskip}
\href{http://www.su.edu}{\includegraphics[height=1.75cm]{sulogo.eps}}
\vspace{-1.69cm}

\begin{center}
  \textbf{\large Math 207: Introduction to Statistics}
\medskip

\textbf{Course Topics}
\bigskip\bigskip
\end{center}

\begin{itemize}
\item Controlled experiments
\item Descriptive statistics
  \begin{itemize}
  \item histograms and bar plots
  \item mean, median, quartiles, percentiles and range
  \item standard deviation and variance
  \item $z$ scores
  \end{itemize}
\item The normal curve
  \begin{itemize}
  \item item
  \end{itemize}
\item Regression
  \begin{itemize}
  \item scatter plots 
  \item correlation $r=\frac{1}{n-1}\,\sum_{i=1}^n\left(\frac{x_i-\mu_x}{\sigma_x}\right)\left(\frac{y_i-\mu_y}{\sigma_y}\right)$
  \item properties of correlation
  \item regression line equation:  $y-\mu_y = r\,\frac{\sigma_y}{\sigma_x}\,(x-\mu_x)$
  \item item
  \item normality of $P(Y\,\vert\,x)$
  \item item
  \end{itemize}
\item Probability
  \begin{itemize}
  \item basic examples with coins, dice, cards and box models
  \item frequency definition of probability and basic properties
  \item conditional probability
  \item multiplication rule and independence
  \item addition rule and mutually exclusive events
  \item binomial formula
  \end{itemize}
\item Sampling distributions
  \begin{itemize}
  \item item
  \end{itemize}
\item Hypothesis testing
  \begin{itemize}
  \item confidence intervals
  \item $z$-test
  \item $t$-test
  \item two-sample $z$-test
  \item $p$-values
  \end{itemize}
\end{itemize}

\vfill
\eject


\begin{center}
  \textbf{\large Math 207: Example Syllabi from Other Institutions}
\bigskip\bigskip
\end{center}


\begin{tabular}{lcl}
Institution & : & University of Illinois at Chicago\\
Course      & : & Stat 101:  Introduction to Statistics\\
%Credits     & : & 3\\
Text        & : & \textit{Intro Stats} by De Veaux, Velleman and Bock, 5th edition\\
\end{tabular}
\bigskip

\begin{tabular}{rll}
1 &	Ch 1--2 & Displaying Data, Bar Chart,Histogram, Descriptive Statistics\\
2 & 	Ch 3--4 &	Contingency Table, Boxplot\\
3 &	Ch 5 	& Normal Distribution,\\
4 & 	Ch 6--7 &	Scatter Plots, Correlation, Linear Regression,Multivariate Regression\\
5 &	Ch 8 	& More about Regression, Exam 1\\
6 &	Ch 9--11 &	Multiple Regression, Surveys, Experiments\\
7 &	Ch 12 	& Probability\\
8 &	Ch 12 	& Probability\\
9 &	Ch 13 	& Sampling Distribution Models, Confidence Intervals Proportions\\
10 &	Ch 14 	& Exam 2\\
11 &	Ch 14 	& CLT, Confidence Intervals for Means\\
12 &	Ch 15 	& Hypothesis testing\\
13 &	Ch 16 	& More about Tests\\
14 &	Ch 17 	& Comparing Two Populations\\
15 &	Ch 18 	& Paired Samples, Review for Final Exam, Final Exam
\end{tabular}

\vfill
\eject
\href{http://www.su.edu}{\includegraphics[height=1.75cm]{sulogo.eps}}
\vspace{-1.69cm}


\begin{center}
  \textbf{\large Math 207: Example Syllabi from Other Institutions}
\bigskip\bigskip
\end{center}


\begin{tabular}{lcl}
Institution & : & Utah State University\\
Course      & : & Stat 1040:  Introduction to Statistics\\
%Credits     & : & 3\\
Text        & : & \textit{Statistics} by Freedman, Pisani and Purves.  4th Ed.
\end{tabular}
\bigskip

\begin{enumerate}
\item Design  of  experiments:  controls,  randomization,  blind  and  double  blind,  placebos.  The  
Salk  vaccine  trial,  historical  controls.  
%
\item Observational   studies:   association   and   causation,   confounding   factors,   the   clofibrate  
trial,  sex-bias  in  graduate  admissions.  
%
\item Descriptive   statistics:   the   histogram,   the   density   scale,   cross-tabulation,   the   average  
and   the   SD   and   their   relationship   to   the   histogram,   the   median,   the   normal  
approximation  for  data,  percentiles,  percentiles  and  the  normal  curve,  measurement  
error,  outliers,  bias  versus  chance  error,  plotting  points  and  lines,  the  slope  and  the  
intercept  of  a  line.  
%
\item Correlation:   the   scatter   diagram,   the   correlation   coefficient,   properties   of   the  
correlation   coefficient   (invariance   to   change   of   location   and   scale,   symmetry),  
ecological  correlations,  correlation  does  not  imply  causation,  examples,  the  SD  line.  
%
\item Regression:  the  graph  of  averages,  regression  to  the  mean,  the  regression  method  for  
individuals,   the   regression   fallacy,   there   are   two   regression   lines,   the   r.m.s.   error,  
plotting   residuals,   heteroscedasticity,   looking   at   vertical   strips,   using   the   normal  
curve  inside  vertical  strips,  the  slope  and  the  intercept,  the  method  of  least  squares,  
does  regression  make  sense?  
%  
\item Probability:  the  long  run  argument,  conditional  probabilities,  the  multiplication  rule,  
independence,   the   Collins   case,   listing   the   ways,   the   addition   rule   the   Paradox   of  
Chevalier  de  Mere,  are  real  dice  fair?  The  binomial  formula.  
%  
\item Chance   variability:   What   does   the   law   of   averages   really   say?   Chance   processes,   the  
sum   of   the   draws,   making   box   models,   the   expected   value   and   the   standard   error,  
using   the   normal   curve,   probability   histograms,   the   normal   approximation   for  
probability  histograms  and  its  scope.  
%
\item Sampling:  The  Literary  Digest  poll,  the  year  the  polls  elected  Dewey,  using  chance  in  
survey  work,  how  well  do  probability  methods  work,  a  close  look  at  the  Gallup  poll,  
telephone  surveys,  chance  error,  bias,  quota  samples,  samples  of  convenience.  
%  
\item Chance   errors   in   sampling:   the   standard   error,   the   correction   factor,   the   Gallup   poll,  
the   accuracy   of   percentages,   confidence   intervals   and   their   interpretation,   the  
accuracy  of  averages,  the  SE  and  expected  value  of  an  average.  
%
\item Tests   of   significance:   null   and   alternative   hypotheses,   test   statistics   and   significance  
levels,  the  role  of  the  box  model,  zero-one  boxes,  the  one-sample  z-tests  and  t-tests,  
the  standard  error  for  a  difference,  comparing  two  sample  averages,  comparing  two  
proportions,   experiments,   the   $\chi^2$-tests   for   independence   and   goodness-of-fit,   how  
Fisher   used   the   $\chi^2$-test   to   show   that   Mendel's   data   was   fudged,   was   the   result  
significant?  Was  it  important?    Data  snooping,  the  importance  of  the  box  model.  
\end{enumerate}

\vfill
\eject
\href{http://www.su.edu}{\includegraphics[height=1.75cm]{sulogo.eps}}
\vspace{-1.69cm}


\begin{center}
  \textbf{\large Math 207: Example Syllabi from Other Institutions}
\bigskip\bigskip
\end{center}


\begin{tabular}{lcl}
Institution & : &  University of Minnesota at Morris\\
Course      & : & Stats 1150:  Introduction to Statistics\\
%Credits     & : & 3\\
Text        & : & \textit{Introduction to the Practice of Statistics} by Moore and McCABE, G. P., 6th Edition
\end{tabular}
\bigskip


{\small\setlength{\baselineskip}{0.97\baselineskip}
 \begin{enumerate}
\item Introduction: Definition and Uses of Statistics
\item Looking at Data: Distributions
   \begin{enumerate}
                  \item Displaying Distributions (excluding Time Plots)
                  \item Describing Distributions
                  \item The Normal Distributions (excluding Assessing Normality, Quantile and Normal Plots)
   \end{enumerate}
\item Looking at Data: Relationships
   \begin{enumerate}
                  \item Scatterplots
                  \item Correlation
                  \item Least-Squares Regression
                  \item Cautions about Regression and Correlation
                  \item The question of Causation
   \end{enumerate}
\item Producing Data
   \begin{enumerate}
                  \item First Steps
                  \item Design of Experiments
                  \item Sampling Design
                  \item Toward Statistical Inference
   \end{enumerate}
\item Probability: The Study of Randomness
   \begin{enumerate}
                  \item Randomness
                  \item Probability Models
                  \item Random Variables
                  \item Mean and Variances of Random Variables
                  \item Probability Laws
   \end{enumerate}
\item Sampling Distributions
   \begin{enumerate}
                  \item Counts and Proportions
                  \item Sample Means
   \end{enumerate}
\item Introduction to Inference
   \begin{enumerate}
                  \item Estimating with Confidence
                  \item Tests of Significance
                  \item Use and Abuse of Tests (excluding Power and Inference as Decision)
   \end{enumerate}
\item Inference for Distributions
   \begin{enumerate}
                  \item Inference for the Mean of a Population
                  \item Comparing Two Means
   \end{enumerate}
\item Inference for Count Data
   \begin{enumerate}
                  \item Inference for a Single Proportion
                  \item Comparing Two Proportions
   \end{enumerate}
\item Inference for Two-Way Tables
   \begin{enumerate}
         \item Data Analysis for Two-Way (Relations in Categorical Data)    
         \item Inference for Two-Way Tables
         \item Formulas and Models for Two-Way Tables  
    \end{enumerate}
\item Inference for Regression
   \begin{enumerate}
                  \item Simple Linear Regression
   \end{enumerate}
\item Multiple Regression
\item One-Way Analysis of Variance
\item Two-Way Analysis of Variance
\item Nonparametric Tests
   \begin{enumerate}
                  \item Wilcoxon Rank Sum Test
                  \item The Wilcoxon Signed Rank Test
                  \item The Kruskal-Wallace Test
   \end{enumerate}
\end{enumerate}}



\vfill
\eject

\href{http://www.su.edu}{\includegraphics[height=1.75cm]{sulogo.eps}}
\vspace{-1.69cm}


\begin{center}
  \textbf{\large Math 207: Example Syllabi from Other Institutions}
\bigskip\bigskip
\end{center}


\begin{tabular}{lcl}
Institution & : & University of Illinois at Chicago\\
Course      & : & Stat 101:  Introduction to Statistics\\
%Credits     & : & 3\\
Text        & : & \textit{Intro Stats} by De Veaux, Velleman and Bock, 5th edition\\
\end{tabular}
\bigskip

%\begin{tabular}{rll}
\vfill
\eject


\end{document}
